\chapter{پروسودی و ویژگی‌های زبرزنجیری گفتار}

در جلسات پیشین، گفتار را عمدتاً از دید «چه گفته شد؟» بررسی کردیم: واج‌ها، فرمنت‌ها، طیف و ضرایب \eng{MFCC}. این جلسه به پرسش مکمل می‌پردازد: «چگونه گفته شد؟». دو گوینده می‌توانند دقیقاً یک دنبالهٔ واجی را بگویند، اما یکی آن را به‌صورت خبر و دیگری به‌صورت پرسش، با اطمینان یا با تردید، سریع یا آهسته ادا کند. این لایهٔ دوم را \term{پروسودی}{Prosody} یا نوای گفتار می‌نامیم. فصل با مروری کوتاه بر جلسهٔ قبل آغاز می‌شود و سپس سه مؤلفهٔ اصلی پروسودی (شدت، فرکانس پایه و زمان)، مشکلات پروسودی در تبدیل متن به گفتار، یک مقالهٔ پژوهشی دربارهٔ اطمینان و صداقت ادراک‌شده و کاربردهای عملی را بررسی می‌کنیم.

%==========================================================
\section{مرور کوتاه جلسهٔ قبل}
%==========================================================

\subsection{مراحل استخراج \eng{MFCC}}
جلسهٔ قبل (فصل پنجم) دیدیم که برای استخراج ضرایب \eng{MFCC} چند گام پشت سر هم روی هر قاب کوتاه از سیگنال انجام می‌شود:
\begin{enumerate}
  \item تبدیل فوریه و محاسبهٔ طیف دامنه یا توان قاب؛
  \item عبور طیف از یک \term{بانک فیلتر مل}{Mel Filterbank} و محاسبهٔ انرژی خروجی هر فیلتر؛
  \item گرفتن لگاریتم از انرژی‌ها؛
  \item اعمال یک تبدیل دیگر برای کاهش هم‌بستگی میان مقادیر؛
  \item نگه‌داشتن تعداد محدودی از ضرایب نخست (مثلاً ۱۲ ضریب) به‌عنوان بردار ویژگی.
\end{enumerate}
فیلترهای مل در فرکانس‌های پایین باریک و در فرکانس‌های بالا پهن‌ترند؛ یعنی جزئیات طیف را در فرکانس‌های پایین دقیق‌تر و در فرکانس‌های بالا «از دورتر و کلی‌تر» می‌بینیم. تشبیه کلاس: مثل عکس‌برداری از مجموعه‌ای از اشیا که بعضی را از نزدیک و با جزئیات و بعضی را از فاصلهٔ دورتر ثبت می‌کنیم.

\begin{correctionbox}{گام چهارم \eng{MFCC} تبدیل فوریهٔ دوباره نیست}
در توضیح شفاهی کلاس از «تبدیل مجدد» (و در جزوه از «تبدیل فوریه مجدد») سخن رفت. در روش استاندارد، این گام \term{تبدیل کسینوسی گسسته}{DCT} است که روی لگاریتم انرژی‌های بانک فیلتر اعمال می‌شود. خروجی \eng{DCT} ضرایبی کم‌هم‌بسته‌تر می‌دهد و بیشتر انرژی اطلاعات در ضرایب نخست متمرکز می‌شود.
\end{correctionbox}

\subsection{پرسش کلاس: عددهای ۴۰ و ۱۲ از کجا می‌آیند؟}
تعداد خروجی‌های بانک فیلتر برابر تعداد فیلترهایی است که طراح انتخاب می‌کند؛ اگر ۴۰ فیلتر بگذاریم، ۴۰ انرژی لگاریتمی خواهیم داشت. اینکه پس از \eng{DCT} چند ضریب نگه داریم (مثلاً ۱۲ یا ۱۳) انتخاب جداگانه‌ای در طراحی ویژگی است. پس تعداد فیلترها و تعداد ضرایب نهایی الزاماً برابر نیستند. نگه‌داشتن ضرایب نخست، نمایش را فشرده می‌کند و عملاً تنها شکل کلی (پوش) طیف را نگه می‌دارد.

\subsection{\eng{MFCC} و مدل منبع–فیلتر}
در مدل \term{منبع–فیلتر}{Source--Filter} تولید گفتار دو بخش دارد: منبع (برای گفتار واک‌دار، ارتعاش تارهای صوتی با فرکانس پایهٔ \eng{$F_0$}) و فیلتر (کل مجرای صوتی). نکتهٔ مورد تأکید کلاس: فیلتر فقط دهان، فک و دندان نیست؛ کل مجرای صوتی از حنجره تا لب را شامل می‌شود. \eng{MFCC} عمدتاً شکل طیفی و ویژگی‌های فیلتر را نگه می‌دارد و اطلاعات دقیق منبع، مانند \eng{$F_0$} و زیر و بمی، در آن تضعیف یا حذف می‌شود. به همین دلیل \eng{MFCC} برای کاربردهایی مثل تشخیص واج و طبقه‌بندی مناسب است (مثلاً ورودی طبقه‌بند یا تصویر ورودی به یک \eng{CNN})، اما برای مطالعهٔ آهنگ جمله کافی نیست. این همان انگیزهٔ ورود به بحث پروسودی است.

نمونهٔ کاربرد مطرح‌شده: طبقه‌بندی نمونه‌های صوتی مرتبط با بیماری (مانند سرفه یا گفتار بیماران کرونا) از نمونه‌های سالم. چون بیماری می‌تواند هم تارهای صوتی و هم مجرای صوتی را تغییر دهد، تغییرات بخش فیلتر ممکن است در ویژگی‌های طیفی دیده شود. این به معنی تشخیص قطعی بیماری با \eng{MFCC} نیست؛ فقط نشان می‌دهد این ویژگی‌ها می‌توانند ورودی مفیدی برای مدل یادگیری ماشین باشند.

\begin{addedbox}{چرا \eng{$F_0$} از \eng{MFCC} حذف می‌شود؟ یک محاسبهٔ عددی}
فرض کنید گفتار با نرخ نمونه‌برداری $16\,\text{kHz}$ ضبط شده و \eng{$F_0=100\,\text{Hz}$} است. دورهٔ تناوب ارتعاش تارها برابر $1/100=10\,\text{ms}$ یعنی $160$ نمونه است. در ظرافت‌های طیف، \eng{$F_0$} به شکل «ساختار هارمونیک» (قله‌هایی با فاصلهٔ $100\,\text{Hz}$) دیده می‌شود. بانک فیلتر مل با پهنای فیلترهای صدها هرتزی این قله‌ها را هموار می‌کند، و نگه‌داشتن فقط ۱۲ ضریب اول \eng{DCT} یعنی نگه‌داشتن تغییرات کند طیف (پوش). بنابراین اطلاعات مربوط به \eng{$F_0$} عملاً در \eng{MFCC} باقی نمی‌ماند.
\end{addedbox}

\subsection{مرور دو واکهٔ «آ» و «اَ»}
پرسش جلسهٔ قبل: تفاوت صدای دو واکهٔ «آ» \eng{/ɑ/} و «اَ» \eng{/æ/} مربوط به منبع است یا فیلتر؟ پاسخ: \textbf{فیلتر}. هر دو واکه واک‌دارند و منبع یکسانی دارند (ارتعاش تارها)؛ اما موقعیت زبان و شکل مجرا فرق می‌کند و در نتیجه فرمنت‌ها جابه‌جا می‌شوند.

\begin{table}[htbp]
\centering
\caption{مقایسهٔ موقعیت زبان و فرمنت‌ها در دو واکهٔ «آ» و «اَ».}
\begin{tabular}{lcc}
\toprule
ویژگی & «آ» \eng{/ɑ/} & «اَ» \eng{/æ/} \\
\midrule
ارتفاع زبان & پایین (باز) & پایین (نسبتاً باز) \\
موقعیت افقی زبان & عقب & جلو \\
\eng{$F_1$} & نسبتاً بالا & نسبتاً بالا \\
\eng{$F_2$} & پایین‌تر & بالاتر \\
\bottomrule
\end{tabular}
\end{table}

دو قاعدهٔ کلیدی: هرچه زبان پایین‌تر باشد \eng{$F_1$} بالاتر است؛ و هرچه زبان جلوتر باشد \eng{$F_2$} بالاتر است. تفاوت اصلی این دو واکه در \eng{$F_2$} است.

\begin{exambox}{نکتهٔ امتحانی}
باید بتوانید با استفاده از \eng{$F_1$} و \eng{$F_2$} دربارهٔ ارتفاع و موقعیت افقی زبان هنگام تولید واکه‌ها توضیح بدهید (مثلاً چرا «آ» و «اَ» با اینکه منبع یکسان دارند متفاوت شنیده می‌شوند).
\end{exambox}

\begin{addedbox}{مقادیر تقریبی فرمنت‌ها}
مقادیر دقیق به گوینده و پیکرهٔ گفتاری بستگی دارد؛ اما برای یک گوینده مرد، به‌طور تقریبی می‌توان انتظار داشت: برای \eng{/ɑ/} حدود \eng{$F_1\approx 650$} تا $750\,\text{Hz}$ و \eng{$F_2\approx 1000$} تا $1200\,\text{Hz}$، و برای \eng{/æ/} حدود \eng{$F_1\approx 600$} تا $750\,\text{Hz}$ و $F_2\approx 1500$ تا $1800\,\text{Hz}$. اختلاف \eng{$F_2$} حدود $500\,\text{Hz}$ یا بیشتر است و همین اختلاف قابل‌شنیدن است، در حالی که \eng{$F_1$} تقریباً مشابه می‌ماند.
\end{addedbox}

%==========================================================
\section{پروسودی چیست؟}
%==========================================================

\subsection{تعریف و ایدهٔ اصلی}
\begin{definitionbox}{پروسودی (نوای گفتار)}
\term{پروسودی}{Prosody} مجموعهٔ ویژگی‌هایی از گفتار است که روی دنبالهٔ آواها و واژه‌ها «سوار» می‌شوند و به آن‌ها آهنگ، ریتم، تکیه و الگوی زمانی می‌دهند. پروسودی نشان می‌دهد \textbf{چگونه} گفته شده است، نه فقط \textbf{چه} گفته شده است.
\end{definitionbox}

مثال: واژهٔ «رفتی» را می‌توان به سه شکل گفت: «رفتی.» (خبری)، «رفتی؟» (پرسشی) و «رفتی!» (تعجب یا تأکید). واج‌ها یکسان‌اند، اما معنا و نیت گوینده فرق می‌کند. علائم نگارشی مانند علامت سؤال یا ویرگول بخشی از این اطلاعات را در متن نشان می‌دهند، ولی تصویر کامل پروسودی در متن نیست؛ برای بررسی دقیق آن باید سیگنال صوتی (موج، اسپکتروگرام و کانتور ویژگی‌ها) را تحلیل کرد.

\subsection{قیاس با عروض در شعر سنتی}
برای نشان دادن «لایه‌ای» بودن پروسودی، استاد به عروض اشاره کرد. در شعر سنتی، علاوه بر واژه‌ها و معنا، شعر روی یک الگوی وزنی سوار شده است و همین الگوست که احساس آهنگ ایجاد می‌کند. علم عروض هجاهای کوتاه (\eng{$\smile$}) و بلند (—) را مشخص می‌کند و از ترتیب آن‌ها وزن شعر را به دست می‌آورد. پس در زبان‌شناسی، همان‌طور که عروض «چگونه گفتن» شعر را بررسی می‌کند، پروسودی نیز علاوه بر «چه گفتن»، «چگونه گفتن» را بررسی می‌کند.

\begin{figure}[htbp]
\centering
\begin{tabular}{c|c|c|c|c}
\toprule
\textbf{هجاها} & «بِ نا مِ» & «خُ دا وَنْ» & «دِ جا نُ» & «خِ رَد» \\
\textbf{خط عروضی} & \eng{$\smile$\;—\;—} & \eng{$\smile$\;—\;—} & \eng{$\smile$\;—\;—} & \eng{$\smile$\;—} \\
\textbf{رکن} & \textbf{فَعولُنْ} & \textbf{فَعولُنْ} & \textbf{فَعولُنْ} & \textbf{فَعَلْ} \\
\bottomrule
\end{tabular}

\bigskip
{\large «خیزید و خز آرید که هنگام خزان است»}\\[3pt]
{\large مفعولُ \quad مفاعیلُ \quad مفاعیلُ \quad فَعولُنْ}
\caption{دو مثال عروضی روی تخته. بالا: تقطیع مصراع «به نام خداوند جان و خرد» با الگوی هجایی \eng{$\smile$ — —} که سه بار و سپس \eng{$\smile$ —} تکرار می‌شود («فعولن فعولن فعولن فعل»). پایین: مصراعی از منوچهری با الگوی «مفعولُ مفاعیلُ مفاعیلُ فعولن». (بازسازی‌شده از تختهٔ کلاس)}
\end{figure}

نکتهٔ مهم این است که برای تحلیل وزن، معنای واژه‌ها اهمیتی ندارد؛ دو عبارت با معنا و واج‌های کاملاً متفاوت می‌توانند وزن یکسانی داشته باشند. یعنی لایهٔ ریتمیک مستقل از هویت واج‌ها و معنا قابل بررسی است. در گفتار عادی نیز وضع همین است: واژه‌ها و واج‌ها یک لایه‌اند و نحوهٔ بیان آن‌ها لایهٔ دیگر.

\begin{addedbox}{تحلیل مصراع‌های تخته}
در مصراع «به نام خداوند جان و خرد»، هجاها به این شکل تقطیع می‌شوند: «بِ‌نا‌م» \eng{$\smile$ — —} ، «خُ‌دا‌وَنْ» \eng{$\smile$ — —} ، «دِ‌جا‌نُ» \eng{$\smile$ — —} ، «خِ‌رَد» \eng{$\smile$ —}. به این وزن (\textbf{متقارب مثمن محذوف}) می‌گویند. همچنین مصراع «خیزید و خز آرید که هنگام خزان است» با الگوی «مفعولُ مفاعیلُ مفاعیلُ فعولن» (هزج مثمن اخرب مکفوف محذوف) سروده شده است. در جلسه، تقطیع مصراع دوم چند بار در گفت‌وگو اصلاح شد و فقط الگوی نهایی روی تخته معتبر است.

دو نکتهٔ تکمیلی: ۱) وزن شعر فارسی \term{کمّی}{Quantitative} است، یعنی بر مبنای کوتاهی و بلندی هجاهاست؛ در مقابل، وزن بسیاری از اشعار انگلیسی بر مبنای تکیه (\eng{stress}) است. ۲) این مثال فقط یک قیاس برای فهم پروسودی است و درس وارد تقطیع شعر نمی‌شود.
\end{addedbox}

\subsection{ویژگی‌های زبرزنجیری}
پروسودی را \term{ویژگی‌های زبرزنجیری}{Suprasegmental Features} نیز می‌نامند. پیشوند \eng{Supra} یعنی «فراتر/بالا» و \eng{Segment} یعنی «قطعه» (واج یا آوا). تصویر کلاس: پروسودی مثل چتری است که از بالا روی قطعات گسترده شده است. دو پیامد مهم:
\begin{itemize}
  \item تحلیل در سطح کل جمله یا عبارت انجام می‌شود، نه فقط یک پنجرهٔ کوتاه یا یک واج. در فصل‌های قبل سیگنال را به پنجره‌های حدود ۲۵ میلی‌ثانیه تقسیم می‌کردیم؛ برای پروسودی باید چند هجا یا یک جمله کامل را با هم دید.
  \item آنچه اهمیت دارد افت و خیزها نسبت به بخش‌های قبل و بعد است، نه مقدار مطلق در یک لحظه. مثلاً برای فهمیدن اینکه روی کدام کلمه تأکید شده، شدت و آهنگ آن کلمه باید با کلمات اطراف مقایسه شود.
\end{itemize}

به همین دلیل، در پروسودی مفهوم \term{کانتور}{Contour} اهمیت بالایی دارد: کانتور یعنی تغییرات یک ویژگی در طول زمان. به‌جای یک مقدار لحظه‌ای از شدت یا فرکانس، مسیر تغییر آن را در کل عبارت بررسی می‌کنیم.

\subsection{سه مؤلفهٔ اصلی}
وقتی از دانشجویان پرسیده شد چه چیزهایی در گفتار پروسودیک است، مواردی مثل پرسشی یا خبری بودن، احساسات، سرعت، مکث و بلندی صدا مطرح شد. در پردازش گفتار باید این مفاهیم را به کمیت‌های قابل اندازه‌گیری تبدیل کنیم؛ مثلاً «عصبانیت» خود یک کمیت صوتی نیست، بلکه باید ببینیم چه چیزهایی در سیگنال تغییر کرده که شنونده آن را عصبانی درک می‌کند. این کمیت‌ها در سه گروه جای می‌گیرند که نقشه‌اش روی تختهٔ کلاس آمده است.

\begin{figure}[htbp]
\centering
\resizebox{\textwidth}{!}{%
\begin{tikzpicture}[rtl, font=\small,
  nd/.style={draw, rounded corners=6pt, align=center, text width=3.1cm, inner sep=5pt, line width=0.6pt},
  note/.style={align=center, text width=5.2cm, font=\footnotesize, text=black!70},
  arr/.style={-{Stealth[length=2.4mm]}, thick, mainblue},
  sk/.style={red!75!black, thick, -{Stealth[length=2mm]}}]
% centre
\node[draw, diamond, aspect=2.2, fill=softred, align=center, inner sep=2pt, line width=0.6pt] (sup) at (0,0) {\lr{suprasegmental}\\\lr{phonetic features}};
\node[note] (supn) at (0,2.9) {\lr{Supra}: فراتر؛ \lr{Segment}: قطعه (بخش). چتری که از بالا روی قطعات گسترده شده است؛ تحلیل در سطح کل جمله یا عبارت، نه فقط یک پنجره یا یک واج؛ افت و خیزها نسبت به بخش قبل و بعد.};
% duration
\node[nd, fill=yellow!30] (dur) at (-6.8,0) {\textbf{\lr{Duration}}\\زمان\\ریتم / مکث\\\lr{punctuation} (کاما)\\\lr{phrasing}\\رفع ابهام نحوی\\تعداد هجا در ثانیه\\کشیده‌شدن};
% intensity
\node[nd, fill=green!20] (int) at (6.3,0) {\textbf{\lr{intensity contour}}\\\lr{energy}\\شدت\\بلندی\\\lr{emphasis}};
\node[note, text width=5cm, font=\small] at (6.3,2.0) {\lr{$\text{Intensity}\propto\text{Magnitude}^{2}$}};
% pitch
\node[nd, fill=cyan!20] (pit) at (0,-3.4) {\textbf{\lr{Pitch Contour}}\\فرکانس\\زیر و بمی\\آهنگ کلام\\\lr{intonation}};
% pitch children
\node[nd, fill=cyan!20, text width=3.4cm] (dec) at (-3.1,-8.2) {کانتور پیچ جمله خبری: زیر و بمی صدا در آخر به سمت بم شدن می‌رود (سقوط فرکانس). نشان‌دهندهٔ قطعیت و پایان است.};
\node[nd, fill=cyan!20, text width=3.4cm] (que) at (0.9,-8.2) {کانتور پیچ جمله سؤالی: بعد از اکسنت، تارهای صوتی کشیده‌تر شده و فرکانس به شدت بالا (زیر) می‌رود («رفتی؟»).};
% intensity children
\node[nd, fill=green!20, text width=3.1cm] (str) at (4.4,-4.6) {\textbf{\lr{stress}}\\یک هجا در درون یک کلمه\\مثل: «مَردی» [اسم مصدر] در برابر «مردی» [یک مرد]};
\node[nd, fill=green!20, text width=3.4cm] (acc) at (8.4,-4.6) {\textbf{\lr{Accent/Focus}}\\یک کلمه در درون یک جمله\\کانتور اینتنسیتی: اکسنت = قله؛ سکوت یا گفتار بی‌تأکید = صاف};
\node[note, text width=4.2cm] (accn) at (8.4,-8.9) {کلمات در فرهنگ لغت و برای لهجهٔ خاص، استرس مشخصی دارند؛ اما این گوینده است که در موقعیت‌های مختلف تصمیم می‌گیرد به کدام کلمه در گفتمان (\lr{Discourse}) اکسنت بدهد.};
% arrows
\draw[arr] (sup.west) -- (dur.east);
\draw[arr] (sup.east) -- (int.west);
\draw[arr] (sup.south) -- (pit.north);
\draw[arr] (pit.south) -- ++(0,-0.4) -| (dec.north);
\draw[arr] (pit.south) -- ++(0,-0.4) -| (que.north);
\draw[arr] (int.south) -- ++(0,-0.4) -| (str.north);
\draw[arr] (int.south) -- ++(0,-0.4) -| (acc.north);
\draw[densely dashed, black!50] (acc.south) -- (accn.north);
% red sketches: fall / rise
\draw[sk] ([xshift=-0.9cm,yshift=-0.2cm]dec.south) -- ++(0.7,-0.8);
\node[red!75!black, font=\footnotesize] at ([xshift=0.8cm,yshift=-0.55cm]dec.south) {افت};
\draw[sk] ([xshift=-0.9cm,yshift=-0.2cm]que.south) to[out=-80,in=180] ++(0.4,-0.7) to[out=0,in=-110] ++(0.8,0.35);
\node[red!75!black, font=\footnotesize] at ([xshift=0.9cm,yshift=-0.55cm]que.south) {خیز};
\end{tikzpicture}}
\caption{نقشهٔ ذهنی ویژگی‌های زبرزنجیری: سه شاخهٔ زمان (\eng{Duration})، زیر و بمی (\eng{Pitch Contour}) و شدت (\eng{Intensity Contour}). شاخهٔ شدت به تکیهٔ واژگانی (\eng{stress}) و تأکید (\eng{Accent/Focus}) تقسیم شده است؛ شاخهٔ زیر و بمی به کانتور جملهٔ خبری و پرسشی. فلش‌های قرمز افت (خبری) و خیز (پرسشی) انتهای کانتور را نشان می‌دهند. (بازسازی‌شده از تختهٔ کلاس)}
\end{figure}

\begin{table}[htbp]
\centering
\caption{سه مؤلفهٔ اصلی پروسودی، کمیت فیزیکی و ادراک متناظر آن‌ها. (جدول افزوده)}
\begin{tabular}{p{2.7cm}p{3.2cm}p{3.2cm}p{4.2cm}}
\toprule
مؤلفه & کمیت فیزیکی & ادراک شنونده & واحد/ابزار رایج \\
\midrule
فرکانس & فرکانس پایه \eng{(\eng{$F_0$})} & زیر و بمی (\eng{Pitch}) & هرتز، نیم‌پرده؛ کانتور \eng{$F_0$} \\
شدت & دامنه، انرژی، شدت & بلندی (\eng{Loudness}) & دسی‌بل؛ کانتور شدت \\
زمان & مدت آوا/هجا، مکث & سرعت، ریتم، کشش & میلی‌ثانیه، هجا در ثانیه \\
\bottomrule
\end{tabular}
\end{table}

%==========================================================
\section{مؤلفهٔ اول: شدت صدا}
%==========================================================

\subsection{واژگان مرتبط}
در این شاخه اصطلاحاتی مثل \eng{Intensity}، \eng{Energy}، \eng{Loudness}، \eng{Amplitude} و \eng{Magnitude} آمده است. استاد تأکید کرد که این‌ها از نظر علمی دقیقاً هم‌معنا نیستند، ولی در این جلسه برای دسته‌بندی کلی در یک گروه قرار گرفته‌اند. رابطهٔ نوشته‌شده روی تخته چنین است:
\[
\text{\lr{Intensity}} \;\propto\; \text{\lr{Magnitude}}^{2}
\]
یعنی شدت (توان در واحد سطح) با مربع دامنه متناسب است. بلندی (\eng{Loudness}) کمیتی \textbf{ادراکی} است و الزاماً با شدت فیزیکی برابر نیست.

\subsection{دسی‌بل}
انسان نسبت‌های شدت را (نه تفاضل‌ها را) حس می‌کند و بازهٔ شدت‌هایی که می‌شنود بسیار گسترده است. به همین دلیل از مقیاس لگاریتمی \term{دسی‌بل}{dB} استفاده می‌شود.

\begin{definitionbox}{تراز شدت و تراز فشار صوت}
\begin{align*}
L_I &= 10\log_{10}\!\left(\frac{I}{I_0}\right)\ \text{dB}, \qquad I_0=10^{-12}\ \text{W/m}^2 \\
L_p &= \text{\lr{SPL}} = 20\log_{10}\!\left(\frac{p}{p_0}\right)\ \text{dB}, \qquad p_0=20\ \mu\text{Pa}
\end{align*}
چون شدت متناسب با مربع فشار (یا دامنه) است ($I\propto p^2$)، ضریب لگاریتم برای دامنه ۲۰ و برای توان یا شدت ۱۰ است. نسبت دو دامنه $A_2/A_1$ معادل اختلاف تراز $20\log_{10}(A_2/A_1)$ دسی‌بل است.
\end{definitionbox}

\begin{addedbox}{مثال‌های عددی دسی‌بل}
\begin{itemize}
  \item دو برابر شدن \textbf{دامنه}: $20\log_{10}2\approx +6.02\,\text{dB}$ (یعنی تقریباً ۶ دسی‌بل). نصف شدن دامنه: $-6\,\text{dB}$.
  \item دو برابر شدن \textbf{توان یا شدت}: $10\log_{10}2\approx +3.01\,\text{dB}$.
  \item ده برابر شدن شدت: $+10\,\text{dB}$؛ ده برابر شدن دامنه (=صد برابر شدن شدت): $+20\,\text{dB}$.
  \item فشار صوت $p=0{,}02\,\text{Pa}$: $L_p = 20\log_{10}\!\frac{0{,}02}{2\times10^{-5}} = 20\log_{10}(1000)=60\,\text{dB SPL}$، که در حد گفت‌وگوی معمولی است.
  \item دو گوینده که هر یک با تراز $60\,\text{dB}$ و مستقل از هم صحبت می‌کنند: توان‌ها جمع می‌شوند و تراز کل $60+10\log_{10}2\approx 63\,\text{dB}$ است، نه ۱۲۰.
  \item در سیگنال دیجیتال، تراز نسبت به بیشینهٔ مقیاس (\eng{dBFS}) سنجیده می‌شود: اگر مقدار مؤثر (\eng{RMS}) سیگنال $0{,}1$ بیشینهٔ مقیاس باشد، $20\log_{10}0{,}1=-20\,\text{dBFS}$.
\end{itemize}
\end{addedbox}

\coursefigure{0.95}{ch06-db-amplitude.png}{سمت چپ: یک تن خالص با دامنه‌های ۱، ۲ و ۱۰. سمت راست: رابطهٔ لگاریتمی دسی‌بل با نسبت دامنه؛ دو برابر شدن دامنه حدود $+6$ دسی‌بل و ده برابر شدن آن $+20$ دسی‌بل است. (شکل افزوده)}

\begin{addedbox}{کانتور شدت چگونه محاسبه می‌شود؟}
برای هر قاب (مثلاً به طول ۲۵ میلی‌ثانیه و گام ۱۰ میلی‌ثانیه) انرژی میانگین $E=\frac1N\sum_{n}x^2[n]$ را حساب می‌کنیم و $10\log_{10}(E+\varepsilon)$ را رسم می‌کنیم ($\varepsilon$ عدد بسیار کوچکی برای جلوگیری از لگاریتم صفر است). دنبالهٔ این مقادیر در طول زمان همان کانتور شدت است؛ نرم‌افزار \eng{Praat} نیز چنین چیزی (با پنجره‌ای نسبتاً بلندتر) نشان می‌دهد. همین مقدار همان ضریب \eng{c0} در \eng{MFCC} را به ذهن می‌آورد که انرژی کل قاب را نشان می‌دهد.
\end{addedbox}

\begin{correctionbox}{بلندی با شدت یکی نیست}
بلندی ادراک‌شده تابع شدت، فرکانس و مدت صوت است. دو صدا با شدت یکسان ولی فرکانس متفاوت بلندی یکسان ندارند؛ گوش در حدود ۲ تا ۵ کیلوهرتز حساس‌ترین است (منحنی‌های هم‌بلندی). قاعدهٔ سرانگشتی روان‌شناسی صوت: افزایش حدود ۱۰ دسی‌بل تقریباً بلندی را دو برابر می‌کند. به همین دلیل گاهی از وزن‌دهی \eng{A} برای نزدیک‌تر کردن اندازه‌گیری به ادراک استفاده می‌شود.
\end{correctionbox}

\subsection{تکیهٔ واژگانی و تأکید}
شاخهٔ شدت روی تخته به دو مفهوم تقسیم شده است که نباید با هم اشتباه شوند.

\begin{definitionbox}{\eng{Stress} در برابر \eng{Accent}}
\textbf{تکیهٔ واژگانی} (\eng{Lexical Stress}): برجستگی یک \textbf{هجا درون یک واژه}. هر واژه در فرهنگ لغت (و برای هر زبان یا لهجهٔ مشخص) الگوی تکیهٔ معینی دارد.

\textbf{تأکید} (\eng{Accent/Focus}): برجسته‌کردن یک \textbf{واژه درون جمله}. این گوینده است که در موقعیت‌های گوناگون تصمیم می‌گیرد به کدام واژه در گفتمان (\eng{Discourse}) تأکید بدهد. در کانتور شدت، تأکید به‌صورت یک \textbf{قله} دیده می‌شود، و بخش‌های سکوت یا گفتار بی‌تأکید تقریباً صاف‌اند.
\end{definitionbox}

\begin{table}[htbp]
\centering
\caption{مقایسهٔ تکیهٔ واژگانی و تأکید.}
\begin{tabular}{p{2.6cm}p{5.2cm}p{5.2cm}}
\toprule
ویژگی & \eng{Stress} (تکیه) & \eng{Accent} (تأکید) \\
\midrule
سطح بررسی & درون واژه (هجا) & درون جمله یا عبارت (واژه) \\
نقش & بخشی از الگوی تلفظ واژه & برجسته‌سازی اطلاعات مورد نظر گوینده \\
وابستگی & به واژه، زبان و لهجه & به قصد گوینده و بافت \\
تغییر در جمله & الگوی پایه معمولاً حفظ می‌شود & می‌تواند بین واژه‌ها جابه‌جا شود \\
\bottomrule
\end{tabular}
\end{table}

مثال‌های کلاس: در «من گفتم که این کار را بکن» تأکید روی «من» است و در «کتاب را بردار» روی «کتاب».

\begin{warningbox}{اشتباه رایج: «اکسنت» در فارسی دو معنا دارد}
واژهٔ \eng{accent} در این درس یعنی «تأکید روی یک واژه در جمله». اما در گفتار روزمره «اکسنت» را برای \textbf{لهجه} (\eng{regional/foreign accent}، مثلاً لهجهٔ شیرازی یا لهجهٔ انگلیسی) هم به کار می‌برند. در این فصل، منظور همان معنای اول است. همچنین واژهٔ \eng{stress} را با «استرس روانی» اشتباه نگیرید؛ در اینجا یعنی «تکیه».
\end{warningbox}

\subsubsection{رابطهٔ \eng{Accent} با \eng{Stress}}
نکته‌ای که استاد چند بار بر آن تأکید کرد: هنگامی که روی یک واژه تأکید می‌کنیم، برجستگی روی همان هجایی قرار می‌گیرد که تکیهٔ واژگانی دارد. یعنی تأکید بیشتر، محل تکیهٔ واژه را عوض نمی‌کند؛ فقط برجستگی هجای تکیه‌دار را تقویت می‌کند. به عبارت تخته: «روی هر کلمه‌ای که تمرکز شود، استرس به قوت خود باقی است و نمی‌توان روی هجایی اکسنت گذاشت که استرس ندارد.» این ایده با \term{شبکهٔ متریک}{Metrical Grid} نمایش داده می‌شود: هر سطح، یک سطح برجستگی است و تأکید می‌تواند فقط به هجایی برسد که در سطح تکیه حضور دارد. استاد گفت این مبحث بیشتر در تحلیل‌های زبان‌شناختی و روان‌شناختی مطرح است و در این درس قرار نیست از آن استفاده تخصصی شود.

\boardfigure{0.8}{ch06-board-pitch-contours.png}{شبکهٔ متریک برای واژه‌های انگلیسی \eng{vitamins} و \eng{California}. هر هجا یک \eng{x} در سطح پایین دارد؛ هجاهای دارای واکهٔ کامل و سپس هجای تکیه‌دار در سطح‌های بالاتر آمده‌اند و تأکید جمله (\eng{accented syll}) فقط می‌تواند روی هجای تکیه‌دار قرار گیرد.}

در این مثال، در \eng{vitamins} هجای نخست و در \eng{California} هجای سوم (\eng{for}) تکیه دارند و تأکید جمله فقط به این هجاها می‌رسد. هجای \eng{Ca} واکهٔ کامل دارد ولی تکیهٔ اصلی ندارد.

\begin{correctionbox}{دقت دربارهٔ تکیه در فارسی}
در زبان فارسی، تکیهٔ واژگانی اغلب قابل‌پیش‌بینی است: در اسم‌ها و صفت‌ها معمولاً بر \textbf{هجای پایانی} (مثلاً «کتاب») قرار می‌گیرد و پسوندهای تصریفی یا ربطی (مانند کسرهٔ اضافه و ضمایر متصل) معمولاً تکیه نمی‌گیرند؛ در فعل‌ها اغلب بر هجای نخست (پیشوند صرفی مانند «می‌» یا «نه‌») است. با این حال، جزئیات (به‌ویژه در مورد پسوندها) در منابع واج‌شناسی محل بحث است. روی تخته برای نشان دادن اینکه محل تکیه می‌تواند معنا را عوض کند، مثال «مَردی» (اسم مصدر، به معنی مردانگی) در برابر «مردی» (یک مرد) آمده است؛ جایگاه دقیق تکیه در این دو صورت را در اینجا قطعی نمی‌کنیم. نکتهٔ روشن این است که در فارسی، تکیه برخلاف انگلیسی کمتر برای تمایز معنایی میان واژه‌ها به کار می‌رود و نقش آهنگ جمله (و نه تکیه) در مثال‌های کلاس مهم‌تر است.
\end{correctionbox}

\begin{correctionbox}{نشانگرهای صوتی تأکید فقط شدت نیستند}
کلاس تأکید را در شاخهٔ «شدت» جای داد و آن را با قله در کانتور شدت شناخت؛ این درست است ولی ناقص. در پژوهش‌های آوایی، نشانگرهای تکیه و تأکید معمولاً ترکیبی‌اند: افزایش \eng{$F_0$} یا تغییر شکل آن، افزایش مدت هجا و کامل‌تر ادا شدن واکه، و در مرتبهٔ بعد افزایش شدت. در بسیاری از آزمایش‌ها شدت ضعیف‌ترین نشانهٔ مستقل از میان این‌هاست. پس در عمل باید هم‌زمان کانتور شدت، \eng{$F_0$} و مدت را بررسی کرد، چنان‌که خود استاد در ادامه به آن اشاره کرد.
\end{correctionbox}

\coursefigure{0.9}{ch06-synthetic-intensity.png}{کانتور شدت برای جملهٔ «کتاب را به علی دادی» در خوانش خنثی (خاکستری) و در حالتی که روی «علی» تأکید شده است (سبز): تأکید به‌صورت قله روی هجای تکیه‌دار آن واژه دیده می‌شود. داده‌ها مصنوعی و فقط برای توضیح مفهوم‌اند. (شکل افزوده)}

\begin{assignmentbox}{تمرین پیشنهادی با \eng{Praat}}
استاد پیشنهاد کرد این بحث را با \eng{Praat} آزمایش کنید: یک واژه (یا جمله) را یک بار عادی و یک بار با تأکید یا لحن عصبانی ضبط کنید و کانتور شدت و کانتور \eng{$F_0$} آن‌ها را مقایسه کنید. بررسی کنید که افزایش شدت و کشش در همان هجای تکیه‌دار رخ می‌دهد یا نه، و محل تکیهٔ واژه تغییر می‌کند یا نه.
\end{assignmentbox}

%==========================================================
\section{مؤلفهٔ دوم: فرکانس و زیر و بمی}
%==========================================================

\subsection{\eng{$F_0$}، \eng{Pitch} و آهنگ}
مفاهیم این شاخه عبارت‌اند از \eng{Pitch}، \eng{Fundamental Frequency}، \eng{Pitch Contour} و \eng{Intonation}. این‌ها دقیقاً هم‌معنا نیستند، اما به یک حوزهٔ مرتبط اشاره دارند.
\begin{itemize}
  \item \term{فرکانس پایه}{\eng{$F_0$}}: کمیتی \textbf{فیزیکی}، یعنی تعداد دورهٔ ارتعاش تارهای صوتی در ثانیه (هرتز).
  \item \term{زیر و بمی}{Pitch}: کمیتی \textbf{ادراکی}؛ آنچه شنونده به‌صورت «زیر» یا «بم» حس می‌کند. در بیشتر موارد با \eng{$F_0$} هم‌ارزی نزدیک دارد و در عمل این دو را یکی می‌گیرند.
  \item \term{کانتور زیر و بمی}{Pitch Contour}: تغییر \eng{$F_0$} در طول یک جمله یا عبارت.
  \item \term{آهنگ جمله}{Intonation}: الگوی منظم و معنادار این تغییرات، که نقش زبانی (مثلاً تمایز خبر و پرسش) دارد.
\end{itemize}

\begin{addedbox}{دامنهٔ \eng{$F_0$} و نیم‌پرده؛ مثال عددی}
\eng{$F_0$} در مردان بزرگسال معمولاً حدود ۸۵ تا ۱۸۰ هرتز و در زنان حدود ۱۶۵ تا ۲۵۵ هرتز است (مقادیر تقریبی). شنوایی انسان زیر و بمی را تقریباً لگاریتمی حس می‌کند، پس برای مقایسهٔ گویندگان و تغییرات، واحد \term{نیم‌پرده}{Semitone} مفیدتر است:
\[
\Delta = 12\log_2\!\frac{f_2}{f_1}\ \ \text{(نیم‌پرده)}.
\]
مثلاً از ۲۰۰ به ۴۰۰ هرتز دقیقاً یک اکتاو (۱۲ نیم‌پرده) است؛ از ۲۰۰ به ۲۵۰ هرتز برابر $12\log_2 1{,}25\approx 3{,}9$ نیم‌پرده؛ و افت \eng{$F_0$} از ۲۰۰ به ۱۵۰ هرتز برابر $12\log_2 0{,}75\approx -5$ نیم‌پرده است. همین ۱۰۰ هرتز تفاضل در ۱۰۰ تا ۲۰۰ هرتز (۱۲ نیم‌پرده) با ۳۰۰ تا ۴۰۰ هرتز (حدود ۵ نیم‌پرده) بسیار فرق دارد.
\end{addedbox}

\begin{addedbox}{\eng{$F_0$} را چگونه از سیگنال درمی‌آوریم؟}
روش‌های رایج (مثل خودهمبستگی در \eng{Praat}) دنبال تناوب سیگنال در پنجره‌های کوتاه می‌گردند. چند نکتهٔ عملی: ۱) \eng{$F_0$} فقط برای قطعات \textbf{واک‌دار} تعریف می‌شود، پس کانتور در بی‌واک‌ها و سکوت‌ها «پاره» است. ۲) خطاهای دو برابر یا نصف شدن (\eng{$F_0$} \eng{octave errors}) رایج‌اند. ۳) برای مقایسهٔ گویندگان، \eng{$F_0$} را معمولاً نسبت به میانگین یا دامنهٔ خود گوینده نرمال می‌کنند (چنان‌که در نمودارهای این فصل نیز دیده می‌شود).
\end{addedbox}

\subsection{کانتور جملهٔ خبری و پرسشی}
دو الگوی اصلی روی تخته:
\begin{itemize}
  \item \textbf{جملهٔ خبری}: زیر و بمی صدا در پایان به سمت بم‌شدن می‌رود (سقوط فرکانس)؛ نشان‌دهندهٔ \textbf{قطعیت و پایان} است.
  \item \textbf{جملهٔ پرسشی}: بعد از اکسنت، تارهای صوتی کشیده‌تر می‌شوند و فرکانس به شدت بالا (زیر) می‌رود (مثال: «رفت؟؟؟»).
\end{itemize}
در جملهٔ پرسشی، افزایش کشش تارهای صوتی (با عمل ماهیچه‌های حنجره) و افزایش فشار هوای زیر تارها باعث بالا رفتن \eng{$F_0$} می‌شود. در شکل زیر همین الگوها روی داده‌های واقعی دیده می‌شود.

\boardfigure{0.95}{ch06-board-metrical-grid.png}{کانتور \eng{$F_0$} نرمال‌شده (برحسب هرتز) در طول زمان نرمال‌شده برای چهار حالت \eng{Statement}، \eng{Question}، \eng{Happy} و \eng{Sad}. جملهٔ خبری در پایان افت می‌کند، پرسشی در پایان بالا می‌رود، حالت شاد در سطح بالاتر و تقریباً هموار است و حالت غمگین ابتدا بالا و سپس با افت شدید تمام می‌شود. دو منحنی رنگی (سیاه و نارنجی) با نوارهای اطمینان خاکستری دو گروه یا شرایط را نشان می‌دهند؛ تعریف آن‌ها در اسلاید ذکر نشده است.}


\begin{correctionbox}{«خبری = افت، پرسشی = خیز» یک گرایش است، نه قانون}
الگوی افت پایانی در جملهٔ خبری و خیز پایانی در پرسش‌های بله/خیر در بسیاری از زبان‌ها (از جمله فارسی) رایج است، اما قاعدهٔ مطلق نیست. پرسش‌های دارای واژهٔ پرسشی (چرا، کجا، چه) اغلب با آهنگ افتان گفته می‌شوند. در فارسی گفتاری، پرسش بله/خیر اغلب بدون «آیا» و فقط با آهنگ ساخته می‌شود («رفتی؟» در برابر «رفتی.»)، و به همین دلیل آهنگ نقش تعیین‌کننده دارد. شکل دقیق کانتور به گوینده، زبان، نوع پرسش و تأکید هم بستگی دارد. در زبان‌های تونال (مثل چینی ماندارین) زیر و بمی در سطح واژه معنای واژگانی دارد؛ فارسی تونال نیست.
\end{correctionbox}

\coursefigure{0.9}{ch06-synthetic-f0.png}{کانتور مصنوعی \eng{$F_0$} برای همان واژه‌ها در دو حالت: جملهٔ خبری با افت پایانی و پرسش بله/خیر با خیز پایانی. هر دو یک «شیب کلی رو به پایین» (\eng{declination}) و قله‌های کوچک روی هجاهای تکیه‌دار دارند. (شکل افزوده)}

\subsection{مثال «کتاب را به علی دادی؟»}
استاد جملهٔ «کتاب را به علی دادی؟» را چند بار با لحن‌های گوناگون گفت. در حالت پرسشی انتظار داریم \eng{$F_0$} در پایان جمله بالا برود؛ در حالت خبری («کتاب را به علی دادی.») زیر و بمی در پایان پایین می‌آید و قطعیت بیشتری منتقل می‌شود. هنگام اجرا، استاد متوجه شد خیز پایانی یکی از دفعات پرسشی به اندازهٔ انتظار روشن نبود. این نشان می‌دهد یک الگوی پروسودیک همیشه یک‌شکل تولید نمی‌شود و در تحلیل باید شمار زیادی نمونه را بررسی کرد.

\subsection{رابطهٔ زیر و بمی با تأکید}
اگر بخواهیم در همین جمله روی «علی» تأکید کنیم، در ناحیهٔ «علی» علاوه بر تغییر شدت، تغییر محسوسی در کانتور \eng{$F_0$} هم دیده می‌شود. بنابراین با بررسی هم‌زمان کانتور شدت و کانتور زیر و بمی می‌توان دربارهٔ محل برجستگی اطلاعات بهتر قضاوت کرد. در کلاس گفته شد پس از بخش تأکیدی، کانتور فرکانس ممکن است به سمت بالا یا پایین حرکت کند؛ پس محل تأکید و ادامهٔ مسیر \eng{$F_0$} به هم مرتبط‌اند.

\begin{addedbox}{اصطلاحات زبان‌شناختی آهنگ}
در زبان‌شناسی، آهنگ را اغلب با چند «برجستگی» و «مرز» توصیف می‌کنند: \term{تکیهٔ زیر و بمی}{Pitch Accent} حرکتی از \eng{$F_0$} روی هجای تکیه‌دار واژهٔ مورد تأکید است، و \term{نغمهٔ مرزی}{Boundary Tone} خیز یا افت در مرز عبارت. روش‌های برچسب‌گذاری مانند \eng{ToBI} بر همین دو مفهوم استوارند. چنین برچسب‌هایی در سامانه‌های \eng{TTS} و پژوهش‌های پروسودی کاربرد دارند.
\end{addedbox}

\subsection{جمله‌های دستوری، مؤدبانه و احساسی}
استاد دربارهٔ جمله‌های دستوری نیز صحبت کرد. در دستور قاطع ممکن است افت پایانی شدیدتر از جملهٔ خبری معمولی باشد؛ اما وقتی جمله حالت درخواست مؤدبانه دارد («بفرمایید، بشینید.») الگوی فرکانس متفاوت است و لحن ملایم‌تر می‌شود. حالت‌های احساسی مثل شادی و غم هم کانتور و سطح متفاوتی دارند (شکل کانتورهای \eng{Happy} و \eng{Sad} بالا). نکتهٔ مهم: الگوی زیر و بمی فقط به کلمات وابسته نیست؛ شیوهٔ بیان، قصد ارتباطی و شرایط گفتار هم در آن مؤثرند.

%==========================================================
\section{مؤلفهٔ سوم: زمان، ریتم و مکث}
%==========================================================

\subsection{ویژگی‌های زمانی}
\begin{table}[htbp]
\centering
\caption{ویژگی‌های زمانی گفتار (شاخهٔ \eng{Duration} در نقشهٔ کلاس).}
\begin{tabular}{p{3.3cm}p{9.5cm}}
\toprule
ویژگی & توضیح \\
\midrule
\eng{Duration} & مدت‌زمان بیان یک آوا، هجا، واژه یا عبارت \\
\eng{Speech Rate} & سرعت گفتار؛ ساده‌ترین معیار «تعداد هجا در ثانیه» است \\
\eng{Rhythm} & الگوی زمانی و ریتمیک گفتار \\
\eng{Pause} & مکث و سکوت \\
\eng{Lengthening} & کشیده‌شدن آواها \\
\eng{Phrasing} & تقسیم گفتار به گروه‌ها و عبارت‌های آهنگین \\
\bottomrule
\end{tabular}
\end{table}

\begin{addedbox}{محاسبهٔ سرعت گفتار؛ مثال عددی}
جملهٔ «مادر کتاب را روی میز اتاق گذاشت و پنجره را باز کرد.» ۱۹ هجا دارد (ما‌دَر، که‌تاب، را، رو‌یِ، میز، او‌تاق، گو‌زاشت، و، پَن‌جِ‌رِ، را، باز، کَرد). فرض کنید ادای آن $4.5$ ثانیه طول بکشد که $0.5$ ثانیه آن مکث است. آنگاه
\[
\text{\lr{speech rate}}=\frac{19}{4{,}5}\approx 4{,}2\ \text{\lr{syll/s}},\qquad
\text{\lr{articulation rate}}=\frac{19}{4{,}0}\approx 4{,}75\ \text{\lr{syll/s}}.
\]
\term{سرعت گفتار}{Speech Rate} مکث‌ها را در زمان کل حساب می‌کند و \term{سرعت تلفظ}{Articulation Rate} فقط زمان گفتار واقعی را. این دو را نباید با هم اشتباه گرفت.
\end{addedbox}

\subsection{نقش مکث و علائم نگارشی}
مکث می‌تواند معنای نحوی یا ارتباطی جمله را تغییر دهد. اگر جمله‌ای را بدون مکث بگوییم، شنونده ممکن است آن را یک عبارت پیوسته بشنود؛ با قرار دادن مکث در نقطه‌ای مشخص، مرز دو عبارت روشن می‌شود. بخشی از علائم نگارشی (ویرگول، دونقطه، نقطه) نشانهٔ محل مکث یا تغییر بخش‌بندی‌اند؛ اما همهٔ علائم زمانی نیستند. علامت سؤال، برای مثال، به ویژگی آهنگی جمله مربوط است.

\subsection{رفع ابهام نحوی}
یک کاربرد مهم ویژگی‌های زمانی، \term{رفع ابهام نحوی}{Syntactic Disambiguation} است. گاهی جمله در متن بیش از یک تفسیر دارد، اما در گفتار، محل مکث‌ها، کشش آواها و گروه‌بندی زمانی به شنونده نشان می‌دهد ساختار مورد نظر چیست.

\begin{addedbox}{مثال: «مردان و زنان باهوش»}
عبارت «مردان و زنان باهوش» دو معنا دارد: «(مردان) و (زنان باهوش)» یا «(مردان و زنان) باهوش». در گفتار، مکث پس از «مردان» تفسیر اول و مکث پس از «زنان» تفسیر دوم را برجسته می‌کند. در شکل زیر، فقط جای مکث عوض شده است. (این مثال در کلاس نبود و برای توضیح مفهوم افزوده شده است.)
\end{addedbox}

\coursefigure{0.9}{ch06-pause-phrasing.png}{دو گروه‌بندی ممکن برای «مردان و زنان باهوش» که فقط در جای مکث تفاوت دارند. مدت‌ها فرضی و برای توضیح‌اند. (شکل افزوده)}

\begin{addedbox}{کشش پایان عبارت}
در بسیاری از زبان‌ها، هجای پایانی یک عبارت کشیده‌تر ادا می‌شود (\eng{phrase-final lengthening}) و گاهی پیش از مرز عبارت مکث می‌آید. بنابراین مدت یک هجا هم از تکیه تأثیر می‌گیرد و هم از جایگاه آن در عبارت؛ در تحلیل پروسودی باید این دو را از هم جدا کرد.
\end{addedbox}

%==========================================================
\section{پروسودی در تبدیل متن به گفتار}
%==========================================================

استاد در ادامه چند نمونه خروجی از یک سامانهٔ فارسی تبدیل متن به گفتار و شبیه‌سازی صدا را پخش کرد تا اهمیت پروسودی در طبیعی بودن گفتار روشن شود.

\subsection{نمونهٔ اول}
متن ورودی: «مادر کتاب را روی میز اتاق گذاشت و پنجره را باز کرد.» در خروجی، مکث و گروه‌بندی کلمات طبیعی نبود. برای مثال، نحوهٔ بیان بخش‌هایی مثل «مادر کتاب» یا «روی میز اتاق» می‌توانست مرزهای عبارتی را درست نشان ندهد. مشکل اصلی این بود که مدل واج‌ها را ساخته بود اما پروسودی مناسب جمله را بازسازی نکرده بود.

\subsection{نمونهٔ دوم}
در نمونهٔ دیگر، جمله‌ای دربارهٔ شرکت \eng{Anthropic} پخش شد. سامانه در به‌کارگیری کسرهٔ اضافه و اتصال کلمات مشکل داشت و در بعضی بخش‌ها کسره‌های نامناسب تولید می‌کرد؛ سرعت گفتار خروجی هم نسبتاً بالا بود. استاد دربارهٔ اینکه داده‌های آموزشی آن سامانه دقیقاً چه بوده اظهار اطمینان نکرد؛ بنابراین از این نمونه‌ها نمی‌توان دربارهٔ شیوهٔ آموزش آن نتیجهٔ قطعی گرفت.

\subsection{چه چیزهایی را یک سامانهٔ \eng{TTS} باید درست تولید کند؟}
کیفیت \eng{TTS} فقط به تلفظ درست واج‌ها وابسته نیست. سامانه باید این‌ها را هم درست تولید کند:
\begin{itemize}
  \item محل مکث‌ها و مرز عبارت‌ها؛
  \item طول هجاها و واژه‌ها؛
  \item ریتم جمله؛
  \item الگوی تأکید؛
  \item اتصال طبیعی واژه‌ها (مانند کسرهٔ اضافه در فارسی)؛
  \item آهنگ و زیر و بمی.
\end{itemize}

\begin{addedbox}{چرا پروسودی در \eng{TTS} دشوار است؟}
متن یک جمله می‌تواند با چندین پروسودی درست خوانده شود؛ یعنی نگاشت متن به پروسودی «یک‌به‌چند» است و چیزهایی مثل قصد گوینده، تأکید و احساس در متن نیامده‌اند. به همین دلیل، سامانه‌های \eng{TTS} امروزی علاوه بر واج‌ها، صراحتاً مدت، \eng{$F_0$} و انرژی هر هجا یا واج را پیش‌بینی (یا از مرجع تقلید) می‌کنند؛ جزئیات معماری در فصل \eng{TTS} آمده است. ارزیابی طبیعی بودن معمولاً با آزمون‌های شنیداری (مثلاً امتیاز میانگین نظر، \eng{MOS}) انجام می‌شود.
\end{addedbox}

%==========================================================
\section{پژوهش: اطمینان و صداقت ادراک‌شده از روی پروسودی}
%==========================================================

\subsection{مسئله و طراحی آزمایش}
در بخش پایانی جلسه، استاد این پرسش را به‌عنوان تمرین طراحی آزمایش مطرح کرد.

\begin{assignmentbox}{تمرین طراحی آزمایش}
آزمایشی طراحی کنید که نشان دهد صداقت و اطمینان یک گوینده، در نحوهٔ بیان او و ویژگی‌های پروسودیک گفتارش، الگوهای مشترکی دارند.
\end{assignmentbox}

نکتهٔ مهم: باید میان دو موضوع فرق گذاشت. (۱) میزان اطمینان یا صداقتی که \textbf{شنونده} از صدای فرد \emph{برداشت} می‌کند؛ (۲) اطمینان یا صداقت \textbf{واقعی} گوینده (راست یا دروغ بودن گفته‌اش). آزمایش ادراکی کلاس مستقیماً موضوع اول را می‌سنجد. وجود یک الگوی پروسودیک، صداقت واقعی فرد را ثابت نمی‌کند.

\subsubsection{فرضیه‌های اولیهٔ کلاس}
از دانشجویان خواسته شد پیش‌بینی کنند گوینده‌ای که مطمئن صحبت می‌کند چه ویژگی‌هایی دارد.
\begin{itemize}
  \item \textbf{شدت}: حدس این بود که شدت بیشتر یا تأکید واضح‌تر احساس تسلط و اطمینان ایجاد می‌کند. اما بلندی صدای افراد یکسان نیست و ممکن است کسی ذاتاً آرام صحبت کند؛ پس مقایسهٔ مقدار مطلق شدت دو فرد درست نیست. استاد پیشنهاد کرد ابتدا ویژگی‌های معمول هر گوینده را به‌عنوان \term{خط پایه}{Baseline} در نظر بگیریم و تغییرات همان فرد را بررسی کنیم.
  \item \textbf{زیر و بمی}: بر اساس بحث کانتور، افت \eng{$F_0$} در پایان عبارت می‌تواند قطعیت را منتقل کند و با اطمینان ادراک‌شده مرتبط باشد.
  \item \textbf{زمان و مکث}: حدس این بود که مکث‌های متعدد در میانهٔ جمله تردید را نشان می‌دهد. یکی از دانشجویان گفت مکث در ابتدا یا انتهای صحبت ممکن است به‌جای تردید، متانت و سنجیدگی را منتقل کند. استاد پذیرفت که \textbf{محل} مکث هم اهمیت دارد، نه فقط تعداد یا جمع زمان آن‌ها.
\end{itemize}
این‌ها فرضیه‌های قابل آزمون‌اند، نه قواعد قطعی دربارهٔ شخصیت یا صداقت افراد.

\subsubsection{طرح آزمایش پیشنهادی}
\begin{center}
\begin{tikzpicture}[rtl, node distance=0.7cm and 0.9cm, >=Stealth,
  box/.style={draw, rounded corners, align=center, minimum width=2.7cm, minimum height=1.3cm, font=\small, fill=blue!6}]
  \node[box] (a) {\RL{۱. جمع‌آوری گفتار}\\\RL{جملهٔ یکسان با}\\\RL{شیوه‌های بیان متفاوت}};
  \node[box, left=of a] (b) {\RL{۲. ارزیابی شنونده}\\\RL{امتیاز اطمینان و}\\\RL{صداقت ادراک‌شده}};
  \node[box, left=of b] (c) {\RL{۳. استخراج ویژگی}\\\lr{Pitch, Intensity,}\\\lr{Duration}};
  \node[box, left=of c] (d) {\RL{۴. تحلیل آماری}\\\RL{هم‌بستگی ویژگی‌ها}\\\RL{با امتیاز شنوندگان}};
  \draw[->] (a) -- (b); \draw[->] (b) -- (c); \draw[->] (c) -- (d);
\end{tikzpicture}
\end{center}

مثلاً از یک گوینده بخواهیم جملهٔ «من مطمئنم که این کار درسته.» را با شیوه‌های مختلف بگوید، نمونه‌ها را برای شنوندگان پخش کنیم و بپرسیم کدام مطمئن‌تر یا صادقانه‌تر به نظر می‌رسد؛ سپس ویژگی‌های صوتی هر نمونه را استخراج کنیم و رابطهٔ آن‌ها با امتیازها را بسنجیم.

\subsection{هم‌بستگی}
\begin{definitionbox}{هم‌بستگی (\eng{Correlation})}
\textbf{هم‌بستگی مثبت}: با افزایش یک متغیر، متغیر دیگر نیز تمایل به افزایش دارد (مثلاً شدت بیشتر همراه با امتیاز اطمینان بیشتر). \textbf{هم‌بستگی منفی}: با افزایش یک متغیر، دیگری تمایل به کاهش دارد (مثلاً مدت طولانی‌تر همراه با امتیاز اطمینان کمتر). نزدیک صفر یعنی رابطهٔ خطی ضعیف است. ضریب هم‌بستگی (پیرسون) بین $-1$ و $+1$ است.
\end{definitionbox}
هم‌بستگی به‌تنهایی رابطهٔ علت و معلولی را ثابت نمی‌کند.

\subsection{مقاله‌ای که بررسی شد}
مقالهٔ نشان‌داده‌شده روی تخته (برگ نخست آن در شکل زیر) با عنوان زیر در سال ۲۰۲۱ در مجلهٔ \eng{Nature Communications} (جلد ۱۲، شمارهٔ مقاله ۸۶۱) منتشر شده است:
\begin{center}
\begin{latin}
\textit{Listeners' perceptions of the certainty and honesty of a speaker\\ are associated with a common prosodic signature}
\end{latin}
\end{center}
نویسندگان: \eng{Louise Goupil}، \eng{Emmanuel Ponsot}، \eng{Daniel Richardson}، \eng{Gabriel Reyes} و \eng{Jean-Julien Aucouturier}.

\boardfigure{0.45}{ch06-paper-title.png}{برگ نخست مقالهٔ مورد بحث (\eng{Goupil et al., Nature Communications, 2021}). چکیده می‌گوید قضاوت شنوندگان دربارهٔ اطمینان و صداقت گوینده بر یک «امضای پروسودیک مشترک» تکیه دارد.}

\begin{addedbox}{آزمایش اصلی مقاله با آنچه روی تخته آمد یکی نیست}
روی تخته، آزمایش به‌صورت «دو نفر جملهٔ «من مطمئنم که این کار درست است» را با پروسودی متفاوت می‌گویند و شنونده انتخاب می‌کند کدام مطمئن‌تر است» ترسیم شده بود. این فقط خلاصهٔ آموزشی ایدهٔ کلی است. در مقاله (مطالعهٔ اول) طبق گزارش‌های منتشرشده: محرک‌ها \textbf{شبه‌واژه}‌های فرانسوی بودند که با الگوریتم‌های پردازش سیگنال، \eng{$F_0$}، بلندی و مدت آن‌ها به‌صورت تصادفی تغییر داده شده بود؛ هر آزمایه شامل دو نسخه از یک شبه‌واژه با پروسودی تصادفی متفاوت بود و شنونده انتخاب می‌کرد کدام نسخه «مطمئن‌تر» (در یک جلسه) یا «صادق‌تر» (در جلسهٔ دیگر) به نظر می‌رسد. از روی انتخاب‌ها، با روش \term{هم‌بستگی معکوس}{Reverse Correlation}، «هستهٔ پروسودیک» هر قضاوت به دست می‌آید: برای هر نقطهٔ زمانی، میانگین مقدار ویژگی در محرک‌های انتخاب‌شده منهای (یا در برابر) محرک‌های ردشده. \eng{$F_0$} و بلندی در ۱۲ نقطه و مدت در ۵ قطعهٔ زمانی نمونه‌برداری شدند. دقیق‌ترین جزئیات (تعداد شرکت‌کنندگان، آزمون‌های آماری و مطالعات بعدی دربارهٔ تأثیر زبان مادری و حافظه) در متن مقاله آمده است و باید از خود مقاله بازبینی شود.
\end{addedbox}

\subsection{خواندن نمودارهای مقاله}
نمودارهای زیر از همان مقاله‌اند. ردیف بالا مربوط به \textbf{اطمینان} (\eng{certainty}، آبی) و ردیف پایین مربوط به \textbf{صداقت} (\eng{honesty}، سبز) است؛ ستون‌ها سه ویژگی \eng{pitch}، \eng{loudness} و \eng{duration} را نشان می‌دهند. محور افقی «قطعه» (\eng{Segment}) است: بازهٔ زمانی هر محرک به ۱۲ قطعهٔ مساوی برای \eng{$F_0$} و بلندی و ۵ قطعه برای مدت تقسیم شده است، یعنی قطعه‌ها \textbf{بخش‌های یک محرک کوتاه} (نه بخش‌های یک جملهٔ بلند) هستند. محور عمودی \eng{Filter Amp (a.u.)} دامنهٔ هستهٔ هم‌بستگی معکوس به واحد دلخواه است؛ خط خطچین صفر یعنی «بدون اثر»، نوار سایه‌دار بازهٔ عدم‌قطعیت، و دایره‌های بزرگ‌تر قطعه‌هایی‌اند که به‌طور معنادار از صفر فاصله دارند (اندازهٔ دایره بر اساس سطح معناداری است؛ استاد در کلاس حدس زده بود اندازهٔ دایره نشان‌دهندهٔ وزنی است و تعریف دقیق آن را نداشت).

\boardfigure{0.95}{ch06-paper-kernels.png}{هسته‌های پروسودیک اطمینان (ردیف بالا) و صداقت (ردیف پایین) برای \eng{pitch}، \eng{loudness} و \eng{duration} در برابر قطعهٔ زمانی. مقدار مثبت یعنی ویژگی در آن قطعه با انتخاب «مطمئن‌تر/صادق‌تر» همراه بوده؛ مقدار منفی یعنی عکس؛ نزدیک صفر یعنی بی‌اثر. (مقاله: \eng{Goupil et al., 2021})}

برداشت کلاس (که با نتیجهٔ مقاله هم‌خوان است):
\begin{table}[htbp]
\centering
\caption{خلاصهٔ الگوهای نمودار.}
\begin{tabular}{p{2.1cm}p{6.1cm}p{5.2cm}}
\toprule
ویژگی & الگوی دیده‌شده & برداشت \\
\midrule
\eng{Pitch} & هسته در ابتدا نزدیک صفر و در نیمهٔ دوم منفی (اطمینان: قطعه‌های ۱۱–۱۲؛ صداقت: قطعه‌های ۶–۹) & فرکانس پایین‌تر در بخش پایانی (افت آهنگ) با اطمینان و صداقت ادراک‌شدهٔ بیشتر همراه است \\
\eng{Loudness} & مثبت در چند قطعه (ابتدا و قطعه‌های ۷–۸ برای اطمینان؛ قطعه‌های ۳–۵ برای صداقت) & شدت بیشتر در این بخش‌ها با اطمینان یا صداقت بیشتر همراه است \\
\eng{Duration} & منفی در قطعه‌های میانی و نزدیک صفر در دو سر & طولانی‌تر بودن (کندتر) با اطمینان و صداقت کمتر همراه است \\
\bottomrule
\end{tabular}
\end{table}

نکات تفسیری استاد: (۱) در پایان جمله، \eng{$F_0$} هم‌بستگی منفی با تأیید صداقت دارد: «فرکانس پایین‌تر، صداقت بیشتر». (۲) تأکید صوتی (بلندی) در میانهٔ جمله با درک قطعیت مرتبط است. (۳) برای مدت، هم‌بستگی «کلاً منفی» است: وقتی مدت یا مکث بیشتر شود، رأی به قطعیت و صداقت کمتر می‌شود («هم‌بستگی معکوس»). (۴) مقدار مثبت در نمودار \eng{$F_0$} یعنی در آن قطعه \eng{$F_0$} بالاتر با انتخاب «مطمئن‌تر» همراه است؛ مقدار منفی یعنی \eng{$F_0$} پایین‌تر؛ و نزدیک صفر یعنی آن تغییر در تصمیم شنونده اثر مشخصی نداشته است. این چهار نکته را می‌توان از شکل بالا تأیید کرد.

\begin{correctionbox}{دقت در تفسیر نمودار}
\begin{itemize}
  \item \textbf{«مکث» یا «مدت»؟} نمودار سوم به مدت قطعه‌ها (در نتیجه سرعت گفتار) مربوط است، نه به «مکث» در معنای سکوت. در مقاله، سرعت گفتار بیشتر (مدت کوتاه‌تر) با اطمینان ادراک‌شدهٔ بیشتر همراه است. حدس کلاس دربارهٔ مکث‌های میانی با نتایج هم‌جهت است، اما برای نتیجه‌گیری دربارهٔ مکث باید ویژگی مکث جداگانه سنجیده شود.
  \item \textbf{«قسمت جمله»}: قطعه‌ها بخش‌های یک محرک کوتاه‌اند؛ تعبیر «ابتدا، میانه و انتهای جمله» در کلاس را باید «ابتدا، میانه و انتهای واژه/محرک» خواند.
  \item \textbf{علامت مقادیر}: در شکل، علامت طوری است که با خلاصهٔ نتیجهٔ مقاله هم‌خوان است: افت \eng{$F_0$} در بخش پایانی، بلندی بیشتر در بخش‌های ابتدایی/میانی و مدت کوتاه‌تر با اطمینان و صداقت بیشتر همراه‌اند. قرارداد علامت (نتیجه منهای مقابل آن) را اگر در سند اصلی لازم است، باید از متن مقاله بازبینی کرد.
  \item \textbf{میانگین یا شکل؟} ایدهٔ مقاله این است که «شکل» کانتور (مثلاً افت یا خیز) مهم است، نه فقط میانگین \eng{$F_0$}. پس عبارت «فرکانس پایین‌تر» را نباید به‌معنای «صدای بم‌تر به‌طور کلی» خواند، بلکه به‌معنای افت زیر و بمی در بخش پایانی است.
  \item \textbf{هم‌بستگی، علیت نیست}؛ و این نتیجه دربارهٔ \textbf{برداشت شنونده} است، نه صداقت واقعی گوینده. الگو مربوط به شرکت‌کنندگان مقاله است و قانون عمومی برای همهٔ افراد و موقعیت‌ها نیست.
\end{itemize}
\end{correctionbox}

\begin{addedbox}{نتیجهٔ کلی مقاله}
طبق چکیدهٔ مقاله، قضاوت دربارهٔ اطمینان و صداقت هر دو بر یک «امضای پروسودیک مشترک» تکیه دارند که مستقل از دانش مفهومی و زبان مادری شنونده است؛ شنوندگان این امضا را به‌طور خودکار استخراج می‌کنند و بر نحوهٔ به‌خاطر سپردن واژه‌ها اثر می‌گذارد. گزارش نویسندگان از الگوی «غیرقابل‌اطمینان»: آهنگ خیزان، شدت کمتر در آغاز، سرعت کمتر، و تغییرپذیری بیشتر \eng{$F_0$} و سرعت. جزئیات را از متن مقاله بازبینی کنید.
\end{addedbox}

%==========================================================
\section{کاربردهای عملی تحلیل پروسودی}
%==========================================================

استاد در خلال بحث مقاله به کاربردهای تحلیل پروسودی در سامانه‌های هوشمند اشاره کرد. یک نمونه ارزیابی نحوهٔ بیان افراد در موقعیت‌هایی مانند \textbf{مصاحبهٔ شغلی} است. چنین سامانه‌ای می‌تواند چند منبع اطلاعات را ترکیب کند:
\begin{enumerate}
  \item \textbf{تحلیل محتوا}: با \eng{ASR}، گفتار به متن تبدیل و تحلیل می‌شود (مثلاً پاسخ به یک پرسش).
  \item \textbf{تحلیل نحوهٔ بیان}: شدت، فرکانس و سرعت گفتار برای تحلیل نشانه‌های مرتبط با اطمینان ادراک‌شده، روانی گفتار یا نحوهٔ ارائهٔ پاسخ.
  \item \textbf{تحلیل چندوجهی}: ترکیب اطلاعات متنی، صوتی و (در بعضی سامانه‌ها) تصویری برای ارائهٔ امتیاز یا برآورد.
\end{enumerate}

\begin{center}
\begin{tikzpicture}[rtl, node distance=0.7cm and 0.8cm, >=Stealth,
  box/.style={draw, rounded corners, align=center, minimum width=2.6cm, minimum height=1cm, font=\small, fill=green!7}]
  \node[box] (s) {\RL{گفتار}\\\RL{مصاحبه‌شونده}};
  \node[box, below left=0.6cm and 1.4cm of s] (asr) {\lr{ASR}\\\RL{محتوای متنی}};
  \node[box, below=0.6cm of s] (pro) {\RL{ویژگی‌های پروسودیک}\\\lr{F0, Intensity, Duration}};
  \node[box, below right=0.6cm and 1.4cm of s] (oth) {\RL{سایر نشانه‌ها}\\\RL{(مثلاً تصویر)}};
  \node[box, below=0.7cm of pro, fill=orange!12] (m) {\RL{مدل ترکیبی}\\\RL{امتیاز / برآورد}};
  \draw[->] (s) -- (asr); \draw[->] (s) -- (pro); \draw[->] (s) -- (oth);
  \draw[->] (asr) -- (m); \draw[->] (pro) -- (m); \draw[->] (oth) -- (m);
\end{tikzpicture}
\end{center}

\begin{warningbox}{احتیاط در استفاده}
چنین امتیازهایی لزوماً بیانگر شخصیت واقعی، تخصص واقعی یا صداقت فرد نیستند. اعتبار آن‌ها نیازمند ارزیابی مستقل و کنترل تفاوت‌های فردی، زبانی و موقعیتی است. پروسودی با خلق، سبک ارتباطی، زبان و فرهنگ و حتی ویژگی‌های فیزیکی گوینده تغییر می‌کند.
\end{warningbox}

\begin{addedbox}{کاربردهای دیگر پروسودی}
علاوه بر ارزیابی گفتار: تشخیص احساس از گفتار، تشخیص نوع جمله و مرز عبارت‌ها (و بازسازی علائم نگارشی در خروجی \eng{ASR})، مدیریت نوبت‌گیری در گفت‌وگو، یادگیری زبان دوم و ارزیابی تلفظ، تشخیص و پایش بیماری‌های گفتاری، و سنتز گفتار طبیعی و گویا.
\end{addedbox}

%==========================================================
\section{جمع‌بندی}
%==========================================================

پروسودی لایه‌ای از گفتار است که «چگونه گفتن» را توصیف می‌کند و تمام اطلاعات آن در متن نوشتاری نیست. سه مؤلفهٔ آن عبارت‌اند از زیر و بمی (\eng{$F_0$} و کانتور آن)، شدت و زمان؛ و چون این ویژگی‌ها زبرزنجیری‌اند، باید آن‌ها را روی بازه‌های بلندتر و به‌صورت کانتور در طول زمان بررسی کرد. شدت در دسی‌بل (با رابطهٔ $20\log_{10}$ برای دامنه) سنجیده می‌شود و در تأکید نقش دارد؛ آهنگ جمله (خبری و پرسشی) عمدتاً در کانتور \eng{$F_0$} دیده می‌شود؛ و ویژگی‌های زمانی در گروه‌بندی و رفع ابهام نحوی مؤثرند. کیفیت \eng{TTS} به پروسودی وابسته است. در پژوهش اطمینان و صداقت ادراک‌شده دیدیم که اثر ویژگی‌ها به جایگاه زمانی آن‌ها بستگی دارد. در نهایت، تحلیل طیفی به‌تنهایی کافی نیست و برای درک نحوهٔ بیان، تأکید و ریتم باید ویژگی‌های پروسودیک را نیز بررسی کرد.

\begin{exambox}{نکات مهم برای مرور و امتحان}
\begin{enumerate}
  \item \eng{MFCC} بیشتر برای نمایش ویژگی‌های طیفی و فیلتر مجرای صوتی است؛ اطلاعات \eng{$F_0$} در آن عملاً حذف می‌شود.
  \item خروجی بانک فیلتر مل به تعداد فیلترها بستگی دارد؛ تعداد ضرایب نهایی (مثلاً ۱۲) می‌تواند کمتر باشد. گام تبدیل پس از لگاریتم \eng{DCT} است.
  \item اختلاف «آ» و «اَ» مربوط به فیلتر است. \eng{$F_1$} با ارتفاع زبان (زبان پایین‌تر، \eng{$F_1$} بالاتر) و \eng{$F_2$} با موقعیت جلو/عقب زبان (جلوتر، \eng{$F_2$} بالاتر) مرتبط است.
  \item پروسودی لایه‌ای از گفتار است که ریتم، آهنگ و تأکید را توصیف می‌کند و تمام آن در متن نیست.
  \item سه مؤلفهٔ پروسودی: \eng{Pitch/Frequency}، \eng{Intensity}، \eng{Time/Duration}.
  \item ویژگی‌های پروسودیک زبرزنجیری‌اند: باید یک عبارت یا دنبالهٔ گفتاری و کانتور آن را بررسی کرد.
  \item \eng{Stress}: تکیه درون واژه؛ \eng{Accent}: برجسته‌سازی واژه در جمله؛ تأکید روی همان هجای تکیه‌دار قرار می‌گیرد.
  \item کانتور شدت می‌تواند ناحیه‌های تأکیدی را نشان دهد، اما \eng{$F_0$} و مدت نیز در تأکید نقش دارند. شدت متناسب با مربع دامنه است؛ $L_p=20\log_{10}(p/p_0)$ و $L_I=10\log_{10}(I/I_0)$.
  \item کانتور جملهٔ پرسشی اغلب با خیز و جملهٔ خبری با افت پایانی همراه است.
  \item ویژگی‌های زمانی: سرعت گفتار، مدت هجا، ریتم، مکث و کشش؛ در \eng{Phrasing} و رفع ابهام نحوی نقش دارند.
  \item برای بررسی پروسودی اطمینان/صداقت: جمع‌آوری گفتار، ارزیابی شنونده، استخراج ویژگی، تحلیل هم‌بستگی.
  \item جایگاه زمانی ویژگی مهم است؛ هم‌بستگی منفی/مثبت/نزدیک صفر را بتوانید در نمودار تفسیر کنید.
\end{enumerate}
\end{exambox}
